\section{Thiết kế AI Assistant}

AI Assistant là một dịch vụ độc lập giúp người vận hành truy vấn và điều khiển
đội thiết bị bằng ngôn ngữ tự nhiên. Nguyên tắc thiết kế cốt lõi là \textbf{Assistant
không bao giờ truy cập trực tiếp cơ sở dữ liệu hay broker}: mọi dữ liệu và thao
tác đều đi qua MCP server của Cloud, nơi cùng dùng các application service với
REST API. Nhờ đó quy tắc nghiệp vụ, kiểm tra dữ liệu và xác thực chỉ được hiện
thực một lần.

\subsection{Công cụ qua Model Context Protocol}

Cloud công bố 17 công cụ qua Model Context Protocol \cite{mcp} (streamable HTTP tại \texttt{/mcp}, bảo vệ bằng
API key), chia làm ba nhóm:

\begin{itemize}
  \item \textbf{Fleet:} \texttt{list\_clusters}, \texttt{get\_cluster},
        \texttt{create\_cluster}, \texttt{list\_nodes}, \texttt{get\_node},
        \texttt{register\_node}, \texttt{list\_runtimes}, \texttt{get\_runtime};
  \item \textbf{Deployment:} \texttt{list\_deployments}, \texttt{get\_deployment},
        \texttt{create\_deployment}, \texttt{update\_deployment},
        \texttt{remove\_deployment}, \texttt{control\_runtime};
  \item \textbf{Monitoring:} \texttt{get\_latest\_telemetry},
        \texttt{get\_telemetry\_history}, \texttt{get\_node\_events}.
\end{itemize}

Schema tham số của công cụ được sinh tự động từ chữ ký hàm Python, và phía
Assistant chuyển mỗi công cụ MCP thành một \texttt{StructuredTool} của LangChain.
Công cụ chỉ được định nghĩa một lần tại Cloud; thêm công cụ mới không cần sửa
Assistant.

\subsection{Kiểm soát thao tác thay đổi}

Các công cụ có tiền tố \texttt{create\_}, \texttt{update\_}, \texttt{remove\_},
\texttt{register\_}, \texttt{control\_} làm thay đổi trạng thái hệ thống. Chúng bị
ẩn hoàn toàn khỏi mô hình trừ khi cấu hình cho phép ghi. Khi được phép, prompt hệ
thống và mô tả của từng công cụ yêu cầu mô hình chỉ gọi sau khi đã mô tả chính
xác hành động và người dùng xác nhận rõ ràng ở một tin nhắn trước đó. Như vậy
có hai lớp bảo vệ: lớp cấu hình (quyết định công cụ có tồn tại với mô hình hay
không) và lớp hội thoại (bắt buộc xác nhận).

\subsection{Đồ thị tác tử}

Một lượt hội thoại được xử lý bằng đồ thị LangGraph \cite{langgraph} bốn nút,
trong đó nút đầu tiên theo mẫu ReAct \cite{yao2023react}
(Hình~\ref{fig:assistant-graph}).

\begin{figure}[H]
\centering
\resizebox{\textwidth}{!}{%
\begin{tikzpicture}[node distance=0.6cm and 0.9cm]
  \node[flowbox, fill=green!6] (start) {Câu hỏi};
  \node[flowbox, right=of start] (react) {\textbf{react}\\gọi tool (ReAct)};
  \node[flowbox, right=of react] (resp) {\textbf{response}\\viết câu trả lời};
  \node[flowbox, right=of resp] (chart) {\textbf{genchart}\\sinh biểu đồ};
  \node[flowbox, right=of chart] (sugg) {\textbf{suggestion}\\gợi ý tiếp};
  \node[flowbox, below=of react, fill=gray!10] (end) {Kết thúc\\(trả lời thẳng)};
  \node[flowbox, above=of react, fill=orange!8] (mcp) {MCP tools\\(Cloud)};
  \draw[flowarrow] (start) -- (react);
  \draw[flowarrow] (react) -- (resp);
  \draw[flowarrow] (resp) -- (chart);
  \draw[flowarrow] (chart) -- (sugg);
  \draw[flowarrow] (react) -- node[right, font=\scriptsize]{không cần tool} (end);
  \draw[flowarrow, <->] (react) -- node[right, font=\scriptsize]{tối đa 6 vòng} (mcp);
\end{tikzpicture}}
\caption{Đồ thị xử lý một lượt hội thoại của AI Assistant}
\label{fig:assistant-graph}
\end{figure}

\begin{itemize}
  \item \textbf{react:} mô hình chính (mặc định Claude Sonnet) lặp tối đa 6 vòng
        \emph{gọi tool $\rightarrow$ nhận kết quả}, chỉ thu thập dữ liệu mà không
        viết câu trả lời. Nếu ngay vòng đầu mô hình không cần tool (chào hỏi, hỏi
        lại khi thiếu đối tượng, hỏi xác nhận), câu trả lời đó được dùng luôn và đồ
        thị kết thúc.
  \item \textbf{response:} tổng hợp kết quả tool thành câu trả lời, phát dần từng
        token về giao diện.
  \item \textbf{genchart:} mô hình nhỏ (Claude Haiku) quyết định dữ liệu có đáng
        vẽ biểu đồ không và sinh đặc tả biểu đồ có cấu trúc; lỗi ở bước này chỉ làm
        mất biểu đồ, không làm hỏng lượt hội thoại.
  \item \textbf{suggestion:} sinh các câu hỏi gợi ý tiếp theo.
\end{itemize}

\textbf{Thực thi tool an toàn.} \texttt{ToolRunner} chạy các lời gọi tool của cùng
một vòng song song (tối đa 8), gộp lời gọi trùng tên và tham số, áp thời gian chờ
60~s cho mỗi tool. Tool không tồn tại, lỗi hoặc quá thời gian được trả về mô hình
dưới dạng dữ liệu lỗi để mô hình tự sửa tham số hoặc giải thích, thay vì làm
dừng cả lượt. Lịch sử hội thoại được giữ theo \texttt{session\_id} bằng checkpointer
của LangGraph, cho phép câu hỏi nối tiếp như ``còn node 2 thì sao?''.

\textbf{Giao thức phát trực tuyến.} Client gửi \texttt{POST /chat/stream} để khởi
tạo lượt và nhận \texttt{stream\_id}, sau đó mở \texttt{GET /chat/stream/\{id\}}
dạng Server-Sent Events. Các sự kiện theo thứ tự: \texttt{status} (đang phân tích,
đang gọi tool nào), \texttt{tool\_end}, \texttt{token}, \texttt{done},
\texttt{chart}, \texttt{suggestions}, hoặc một sự kiện \texttt{error} duy nhất.

\subsection{So sánh với hướng tiếp cận RAG}

Trong bài toán vận hành, câu hỏi chủ yếu hướng tới \emph{trạng thái hiện tại}
của hệ thống (node nào offline, runtime nào lỗi, FPS vừa giảm bao nhiêu). Dữ
liệu này thay đổi liên tục nên truy xuất qua tool có cấu trúc chính xác hơn so
với tìm kiếm văn bản đã lập chỉ mục \cite{lewis2020rag}. Vì vậy đồ án ưu tiên
tool calling qua MCP; RAG trên tài liệu, runbook được giữ ở vai trò mở rộng cho
các câu hỏi về quy trình và kiến thức ổn định.
